\documentclass[11pt]{article}

\usepackage[margin=1in]{geometry}
\usepackage{amsmath,amsthm,amssymb}

\newtheorem{theorem}{Theorem}

\begin{document}

\begin{theorem}[Factor bound from the half-shifted ratio]
Let \(A,B,x\in\mathbb{R}\) and \(a\in(0,1/2]\). Assume
\[
    0 < B < A,
    \qquad
    0 \le x < B.
\]
Then
\[
    \frac{A+(2a-1)x}{B+(2a-1)x}
    \le
    \frac{1+2a}{4a}
    \frac{A+\frac12(2a-1)x}{B+\frac12(2a-1)x}.
\]
Consequently, if
\[
    \frac{A+\frac12(2a-1)x}{B+\frac12(2a-1)x}\le U,
\]
then
\[
    \frac{A+(2a-1)x}{B+(2a-1)x}
    \le
    \frac{1+2a}{4a}U.
\]
\end{theorem}

\begin{proof}
Set
\[
    t := (2a-1)x.
\]
Since \(a\in(0,1/2]\), we have \(2a-1\le 0\). Since \(x\ge 0\), it follows that
\[
    t\le 0.
\]
It is convenient to write
\[
    s := -t = (1-2a)x.
\]
Then \(s\ge 0\), and the two ratios become
\[
    \frac{A+t}{B+t}
    =
    \frac{A-s}{B-s},
    \qquad
    \frac{A+\frac12 t}{B+\frac12 t}
    =
    \frac{A-\frac12 s}{B-\frac12 s}.
\]
Because \(x<B\), we have
\[
    s=(1-2a)x < (1-2a)B.
\]
Therefore
\[
    B-s > B-(1-2a)B = 2aB > 0,
\]
and also
\[
    B-\frac12s > 0.
\]
Moreover, since \(s<B<A\), both \(A-s\) and \(A-\frac12s\) are positive.
Thus all denominators below are positive, and the ratios below are positive.

We compare the full-shifted ratio to the half-shifted ratio:
\[
    \frac{
        \frac{A-s}{B-s}
    }{
        \frac{A-\frac12 s}{B-\frac12 s}
    }
    =
    \frac{A-s}{A-\frac12 s}
    \cdot
    \frac{B-\frac12 s}{B-s}.
\]
Since \(s\ge 0\), we have
\[
    A-s \le A-\frac12s,
\]
so
\[
    \frac{A-s}{A-\frac12s}\le 1.
\]
Hence
\[
    \frac{
        \frac{A-s}{B-s}
    }{
        \frac{A-\frac12 s}{B-\frac12 s}
    }
    \le
    \frac{B-\frac12s}{B-s}.
\]
Using the bound \(s < (1-2a)B\), we estimate
\[
    B-\frac12s
    \le B,
\]
and
\[
    B-s > 2aB.
\]
This already gives the coarser bound
\[
    \frac{B-\frac12s}{B-s} < \frac{1}{2a}.
\]
For the sharper stated constant, use the monotonicity of the function
\[
    g(s) := \frac{B-\frac12s}{B-s}
\]
on \(0\le s<B\). Indeed,
\[
    g'(s)
    =
    \frac{\frac12B}{(B-s)^2}
    >0.
\]
Thus \(g\) is increasing, and since \(s < (1-2a)B\),
\[
    g(s)
    <
    g\bigl((1-2a)B\bigr).
\]
Computing the endpoint value gives
\[
    g\bigl((1-2a)B\bigr)
    =
    \frac{B-\frac12(1-2a)B}{B-(1-2a)B}
    =
    \frac{\frac12(1+2a)B}{2aB}
    =
    \frac{1+2a}{4a}.
\]
Therefore
\[
    \frac{
        \frac{A-s}{B-s}
    }{
        \frac{A-\frac12 s}{B-\frac12 s}
    }
    \le
    \frac{1+2a}{4a}.
\]
Multiplying by the positive half-shifted ratio yields
\[
    \frac{A-s}{B-s}
    \le
    \frac{1+2a}{4a}
    \frac{A-\frac12 s}{B-\frac12 s}.
\]
Substituting \(s=-(2a-1)x\) proves
\[
    \frac{A+(2a-1)x}{B+(2a-1)x}
    \le
    \frac{1+2a}{4a}
    \frac{A+\frac12(2a-1)x}{B+\frac12(2a-1)x}.
\]
The final implication follows immediately: if the half-shifted ratio is at most \(U\), then the full-shifted ratio is at most
\[
    \frac{1+2a}{4a}U.
\]
\end{proof}

\end{document}
